\documentclass[11pt]{amsart}
\usepackage{geometry}                % See geometry.pdf to learn the layout options. There are lots.
\geometry{letterpaper}                   % ... or a4paper or a5paper or ... 
%\geometry{landscape}                % Activate for for rotated page geometry
%\usepackage[parfill]{parskip}    % Activate to begin paragraphs with an empty line rather than an indent
\usepackage{graphicx}
\usepackage{amssymb}
\usepackage{epstopdf}
\DeclareGraphicsRule{.tif}{png}{.png}{`convert #1 `dirname #1`/`basename #1 .tif`.png}

\title{Problema 4.17}

%\date{}                                           % Activate to display a given date or no date

\begin{document}
\maketitle
%\section{}
%\subsection{}
\noindent Dada una matriz A(M x N) de tipo entero, construya un diagrama de flujo para
calcular la traspuesta de dicha matriz. La traspuesta de una matriz se obtiene al
escribir las filas de la matriz A como columnas. Por ejemplo si tenemos la siguiente matriz A:
\newline

\[A
\begin{bmatrix}
-5 & 6 & 8 & 4 \\
2 & 5 & 7 & 9 \\
1 & 3 & 4 & 6 \\
\end{bmatrix}
\]
\newline

La traspuesta de dicha matriz es la siguiente:
\newline

\[A^{T}
\begin{bmatrix}
-5 & 6 & 8 \\
4 & 2 & 1 \\
5 & 3 & 7 \\
9 & 6 & 4 \\
\end{bmatrix}
\]
\newline

\noindent Dato: A[1..M, 1..N], 1 $\leq$ M $\leq$ 50, 1 $\leq$ N $\leq$ 30
\newline

\noindent Donde: A es un arreglo bidimensional de tipo entero.
\newline


\noindent\textbf{Explicación de las variables}
\newline

\noindent  M, N: Variables de tipo entero. Indican el número de renglones y de columnas.

\noindent  I,J: Variables de tipo entero. Se usan para el control de ciclos y como índices de los arreglos.

\noindent  A: Arreglo bidimensional de tipo entero.

\noindent  TA: Arreglo bidimensional de tipo entero, TA[1..N, 1..M]. Almacena la traspuesta de la matriz dada.


\end{document}  